% !TeX root = ../../infdesc.tex
\section{Propositional logic}
\secbegin{secPropositionalLogic}

\begin{exercise}
\label{exJohnMathematicianPittsburgh}
Express the proposition `John is a mathematician who lives in Pittsburgh' in the form $p \wedge q$, for propositions $p$ and $q$.
\end{exercise}

\begin{solutions*}
    Let proposition $p$ be `John is a mathematician';

    Let proposition $q$ be `John lives in Pittsburgh'.
    
    Then, the proposition `John is a mathematician who lives in Pittsburgh' can be expressed as $p \wedge q$.
\end{solutions*}

\begin{exercise}
Write a proof tree whose conclusion is the propositional formula $(p \wedge q) \wedge (r \wedge s)$, where $p,q,r,s$ are propositional variables. Express `$2$ is an even prime number and $3$ is an odd prime number' in the form $(p \wedge q) \wedge (r \wedge s)$, for appropriate propositions $p$, $q$, $r$ and $s$, and describe how your proof tree suggests what a proof might look like.
\end{exercise}

\begin{solutions*}
     \fixlistskip
\begin{center}
    \begin{prooftree}
        \AxiomC{$p$}
        \AxiomC{$q$}
    \BinaryInfC{$p \wedge q$}
        \AxiomC{$r$}
        \AxiomC{$s$}
    \BinaryInfC{$r \wedge s$}
    \BinaryInfC{$(p \wedge q) \wedge (r \wedge s)$}
    \end{prooftree}
\end{center}
Let $p$ be `$2$ is an even number'; $q$ be `$2$ is a prime number';

Let $r$ be `$3$ is an odd number'; $s$ be `$3$ is a prime number';

Then the proposition `$2$ is an even prime number and $3$ is an odd prime number' can be represented in the form $(p \wedge q) \wedge (r \wedge s)$. 

When we try to prove this proposition, according to the definition of conjunction, we can firstly prove that $(p \wedge q)$ is true, and then prove that $(r \wedge s)$ is true. One step deeper, we need to prove that $p$ is true, $q$ is true, $r$ is true, and $s$ is true.

\end{solutions*}

\begin{exercise}
Let $n$ be an integer. Prove that $n^2$ leaves a remainder of $0$, $1$ or $4$ when divided by $5$.
\end{exercise}

\begin{solutions*}
    This problem is similar to the Proposition 1.1.18. 

    \begin{proof}
    Let $n \in \mathbb{Z}$. By the division theorem, one of the following must be true for some $k \in \mathbb{Z}$:
    \[
    n=5k \quad \text{or} \quad n=5k+1 \quad \text{or} \quad n=5k+2 \quad \text{or} \quad n=5k+3 \quad \text{or} \quad n=5k+4
    \]
    \begin{itemize}
    \item Suppose $n=5k$. Then
    \[
    n^2 = (5k)^2 = 25k^2 = 5 \cdot (5k^2)
    \]
    So $n^2$ leaves a remainder of $0$ when divided by $5$.
    \item Suppose $n=5k+1$. Then
    \[
    n^2 = (5k+1)^2 = 25k^2+10k+1 = 5(5k^2+2k)+1
    \]
    So $n^2$ leaves a remainder of $1$ when divided by $5$.
    \item Suppose $n=5k+2$. Then
    \[
    n^2 = (5k+2)^2 = 25k^2+20k+4 = 5(5k^2+4k)+4
    \]
    So $n^2$ leaves a remainder of $4$ when divided by $5$.
    \item Suppose $n=5k+3$. Then
    \[
    n^2 = (5k+3)^2 = 25k^2+30k+9 = 5(5k^2+6k+1)+4
    \]
    So $n^2$ leaves a remainder of $4$ when divided by $5$.
    \item Suppose $n=5k+4$. Then
    \[
    n^2 = (5k+4)^2 = 25k^2+40k+16 = 5(5k^2+8k+3)+1
    \]
    So $n^2$ leaves a remainder of $1$ when divided by $5$.
    \end{itemize}
    In all cases, $n^2$ leaves a remainder of $0$ or $1$ or $4$ when divided by $5$.
    \end{proof}
\end{solutions*}

\begin{exercise}
Let $a,b \in \mathbb{R}$ and suppose $a^2-4b \ne 0$. Let $\alpha$ and $\beta$ be the (distinct) roots of the polynomial $x^2+ax+b$. Prove that there is a real number $c$ such that either $\alpha-\beta = c$ or $\alpha - \beta = ci$.
\end{exercise}

\begin{solutions*}
    \fixlistskip
    \begin{proof}
    By the Quadratic formula, we have the roots of the polynomial $x^2+ax+b$ as
    \[
    \frac{-a \pm \sqrt{a^2-4b}}{2}
    \]
    Without loss of generality, let $\alpha = \dfrac{-a+\sqrt{a^2-4b}}{2}$ and $\beta = \dfrac{-a-\sqrt{a^2-4b}}{2}$. Then
    \[
    \alpha-\beta = \dfrac{a^2-4b}{2} = \dfrac{a^2-4b}{2} = \sqrt{a^2-4b}
    \]
    \begin{itemize}
        \item If $a^2-4b > 0$, then $\sqrt{a^2-4b}$ is a real number. In this case, $\exists c \in \mathbb{R}$ such that $\alpha-\beta = c$.
        \item If $a^2-4b < 0$, then $\sqrt{a^2-4b}$ is a complex number. In this case, $\exists c \in \mathbb{R}$ such that $\alpha - \beta = ci$.
    \end{itemize}
    In all cases, we proved that exist a real number $c$ such that either $\alpha-\beta = c$ or $\alpha - \beta = ci$.  
    \end{proof}
\end{solutions*}


\begin{exercise}
Let $p(x)$ be a polynomial over $\mathbb{C}$. Prove that if $\alpha$ is a root of $p(x)$, and $a \in \mathbb{C}$, then $\alpha$ is a root of $(x-a)p(x)$.
\end{exercise}

\begin{solutions*}
    \fixlistskip
    \begin{proof}
    Let $p(x)$ be a polynomial over $\mathbb{C}$. Suppose $\alpha$ is a root of $p(x)$, and $a \in \mathbb{C}$. Then
    \[
    p(\alpha) = 0
    \]
    Substitute $\alpha$ into $(x-a)p(x)$, we get
    \[
    (x-a)p(x) = (\alpha-a)p(\alpha) = 0
    \]
    This proves that $\alpha$ is a root of $(x-a)p(x)$.
    \end{proof}
\end{solutions*}

\begin{exercise}
\label{exQuadraticFormulaConverse}
Complete the proof of the quadratic formula. That is, for fixed $a,b \in \mathbb{C}$, prove that if
\[
\alpha = \frac{-a+\sqrt{a^2-4b}}{2} \quad \text{or} \quad \alpha =\frac{-a-\sqrt{a^2-4b}}{2}
\]
then $\alpha$ is a root of the polynomial $x^2+ax+b$.
\end{exercise}

\begin{solutions*}
    \fixlistskip
    \begin{proof}
    Next we prove that if $\alpha$ is one of the values given in the statement of the proposition, then $\alpha$ is a root of $x^2+ax+b$.
    \begin{itemize}
        \item If $\alpha = \dfrac{-a+\sqrt{a^2-4b}}{2}$, substitute $\alpha$ into $x^2+ax+b$, we get
        \begin{align*}
        x^2+ax+b &= \left(\frac{-a+\sqrt{a^2-4b}}{2}\right)^2 + a\left(\frac{-a+\sqrt{a^2-4b}}{2}\right) + b \\
        &= \frac{a^2-2a\sqrt{a^2-4b}+a^2-4b}{4} + \frac{-a^2+a\sqrt{a^2-4b}}{2} + b \\
        &= \frac{a^2-a\sqrt{a^2-4b}}{2}-b + \frac{-a^2+a\sqrt{a^2-4b}}{2} + b \\
        &= 0
        \end{align*}
        This proves that $\dfrac{-a+\sqrt{a^2-4b}}{2}$ is a root of $x^2+ax+b$.
        \item If $\alpha = \dfrac{-a-\sqrt{a^2-4b}}{2}$, substitute $\alpha$ into $x^2+ax+b$, we get
        \begin{align*}
        x^2+ax+b &= \left(\frac{-a-\sqrt{a^2-4b}}{2}\right)^2 + a\left(\frac{-a-\sqrt{a^2-4b}}{2}\right) + b \\
        &= \frac{a^2+2a\sqrt{a^2-4b}+a^2-4b}{4} + \frac{-a^2-a\sqrt{a^2-4b}}{2} + b \\
        &= \frac{a^2+a\sqrt{a^2-4b}}{2}-b + \frac{-a^2-a\sqrt{a^2-4b}}{2} + b \\
        &= 0
        \end{align*}
        This proves that $\dfrac{-a-\sqrt{a^2-4b}}{2}$ is a root of $x^2+ax+b$.
    \end{itemize}
    In all cases, we proved that $\alpha$ is a root of $x^2+ax+b$.
    \end{proof}
\end{solutions*}

\begin{exercise}
Prove that a natural number $n$ is divisible by $3$ if and only if the sum of its base-$10$ digits is divisible by $3$.
\end{exercise}

\begin{solutions*}
    \fixlistskip
    \begin{proof}
    The base-$10$ expansion of a natural number $n$ is 
    \[
    n = d_r \cdot 10^r + d_{r-1} \cdot 10^{r-1} + \cdots + d_1 \cdot 10 + d_0
    \]
    We can write $10^r$ as $(10^r - 1) + 1$, where $10^r - 1$ is divisible by $3$.
    \begin{align*}
        n &= d_r \cdot (10^r - 1 + 1) + d_{r-1} \cdot (10^{r-1}-1 + 1) + \cdots + d_1 \cdot (10 - 1 + 1) + d_0\\
        &= d_r \cdot (10^r - 1) + d_r + d_{r-1} \cdot (10^{r-1}-1) + d_{r-1} + \cdots + d_1 \cdot (10-1) + d1 + d_0 \\
        &= d_r \cdot (10^r - 1) + d_{r-1} \cdot (10^{r-1}-1) + \cdots + d_1 \cdot (10-1) + (d_r + d_{r-1} + \cdots + d_1 + d_0) \\
    \end{align*}
    Since $3$ divides $10^r - 1$, then 
    \[n \equiv (d_r + d_{r-1} + \cdots + d_1 + d_0) \pmod{3}\]
    \begin{itemize}
        \item ($\implies$) Suppose $3$ divides $n$. It follows that $3$ divides $(d_r + d_{r-1} + \cdots + d_1 + d_0)$.
        \item ($\impliedby$) Suppose $3$ divides $(d_r + d_{r-1} + \cdots + d_1 + d_0)$. It follows that $3$ divides $n$.
    \end{itemize}
    This proves that a natural number $n$ is divisible by $3$ if and only if the sum of its base-$10$ digits is divisible by $3$.
    \end{proof}
\end{solutions*}

\begin{exercise}
\label{exNegationAndReciprocalOfIrrationalNumbers}
Let $x \in \mathbb{R}$. Prove by contradiction that if $x$ is irrational then $-x$ and $\frac{1}{x}$ are irrational.
\end{exercise}

\begin{solutions*}
    \fixlistskip
    \begin{proof}
    Assume for the sake of contradiction that $x$ is irrational, $-x$ and $\frac{1}{x}$ are rational. Then we can represent $-x$ as $-x = \frac{p}{q}$ for some integers $p$ and $q$ with $q \ne 0$. 
    \[-x = \frac{p}{q} \implies x = -\frac{p}{q}\]
    This contradicts the assumption that $x$ is irrational. It follows that $-x$ is not rational, and is therefore irrational.

    Similarly, we can represent $\frac{1}{x}$ as $\frac{1}{x} = \frac{s}{r}$ for some integers $s$ and $r$ with $r \ne 0$. 
    \[\frac{1}{x} = \frac{s}{r} \implies x = \frac{r}{s}\]
    This contradicts the assumption that $x$ is irrational. It follows that $\frac{1}{x}$ is not rational, and is therefore irrational.
    \end{proof}
\end{solutions*}

\begin{exercise}
\label{exNoLeastPositiveReal}
Prove by contradiction that there is no least positive real number. That is, prove that there is not a positive real number $a$ such that $a \le b$ for all positive real numbers $b$.
\end{exercise}

\begin{solutions*}
    \fixlistskip
    \begin{proof}
    Assume for the sake of contradiction that there exists a positive real number $a$ such that $a \le b$ for all positive real numbers $b$. In other words, there is no positive real number less than $a$.

    We can easily construct a counterexample.

    Let $a' = \dfrac{a}{2}$. $a'$ is positive, and $a' < a$. 

    This contradicts the assumption that $a \le b$ for all positive real numbers $b$. It follows that there is no least positive real number.
    \end{proof}
\end{solutions*}

\begin{exercise}
Reflect on the proof of Proposition 1.1.46. Where in the proof did we use the law of excluded middle? Where in the proof did we use proof by contradiction? What was the contradiction in this case? Prove Proposition 1.1.46 twice more, once using contradiction and not using the law of excluded middle, and once using the law of excluded middle and not using contradiction.
\end{exercise}

\begin{solutions*}
\fixlistskip
\begin{itemize}
    \item \textbf{Where was the Law of Excluded Middle used?}\\
        The Law of Excluded Middle was used at the very beginning to establish the case structure: "Suppose $a$ is even... Suppose $a$ is odd." This implicitly relies on the axiom that $P \lor \neg P$ is true.
    \item \textbf{Where was Proof by Contradiction used?}\\
        It was used entirely within the second case. Instead of proving directly that $b$ must be even when $a$ is odd, the proof assumed the opposite ("If $b$ is also odd...") to see if it would break the universe.
    \item \textbf{What was the contradiction?}\\
        The contradiction was the simultaneous truth of two mutually exclusive statements: the premise established that $ab$ is even, but the logical deduction from assuming $b$ was odd resulted in $ab$ being odd.
\end{itemize}
 
\textbf{Proof 1: Using Contradiction (Without LEM)}

To avoid the Law of Excluded Middle, we do not start by splitting $a$ into "even or odd." Instead, we assume the negation of the entire conclusion and seek a contradiction.
\begin{proof}
    Suppose $ab$ is even. Assume for the sake of contradiction that ``either $a$ or $b$ is even'' is false, that is $a$ and $b$ are both odd.

    By Proposition 1.1.42, we have $a = 2k+1$ and $b = 2\ell+1$ for some $k, \ell \in \mathbb{Z}$. Then
    \[ab = (2k+1)(2\ell+1) = 4k\ell + 2k + 2\ell + 1 = 2(2k\ell + k + \ell) + 1\]
    Since $(2k\ell + k + \ell) \in \mathbb{Z}$, by Proposition 1.1.42, $ab$ is odd. This contradicts the assumption that $ab$ is even. So $a$ and $b$ cannot both be odd. Therefore, if $ab$ is even, then either $a$ is even or $b$ is even (or both).
\end{proof}

\textbf{Proof 2: Using LEM (Without Contradiction)}

To avoid the Contradiction, we use LEM to set up cases and then prove each case directly without assuming the opposite of what we want to find.
\begin{proof}
    By LEM, $a$ is either even or odd.
    \begin{itemize}
        \item \textbf{Case 1: $a$ is even.}\\
          The statement "$a$ is even or $b$ is even" is immediately true.
        \item \textbf{Case 2: $a$ is odd.}\\
          We must now show $b$ is even. By LEM, $b$ is either even or odd.
          \begin{itemize}
            \item \textbf{Case 2.1: $b$ is even.}\\
                the statement is true.
            \item \textbf{Case 2.2: $b$ is odd.}\\
                By Proposition 1.1.42, we have $a = 2k+1$ and $b = 2\ell+1$ for some $k, \ell \in \mathbb{Z}$. Then
                \[ab = (2k+1)(2\ell+1) = 4k\ell + 2k + 2\ell + 1 = 2(2k\ell + k + \ell) + 1\]
                Since $(2k\ell + k + \ell) \in \mathbb{Z}$, by Proposition 1.1.42, $ab$ is odd. 
          \end{itemize}
    \end{itemize}
    In summary, either $a$ is even or $b$ is even (or both), $ab$ is even.
\end{proof}
\end{solutions*}

\begin{exercise}
Let $a$ and $b$ be irrational numbers. By considering the number $\sqrt{2}^{\sqrt{2}}$, prove that it is possible that $a^b$ be rational.
\end{exercise}

\begin{solutions*}
    \fixlistskip
    \begin{proof}
    We know that $\sqrt{2}$ is an irrational number. Consider the number $\sqrt{2}^{\sqrt{2}}$. There are two possible cases:
    \begin{itemize}
    \item If $\sqrt{2}^{\sqrt{2}}$ is rational, then we can choose $a = \sqrt{2}$ and $b = \sqrt{2}$, and $a^b = \sqrt{2}^{\sqrt{2}}$ is rational.
    \item If $\sqrt{2}^{\sqrt{2}}$ is irrational, then we can choose $a = \sqrt{2}^{\sqrt{2}}$ and $b = \sqrt{2}$,
    \[
    a^b = \left(\sqrt{2}^{\sqrt{2}}\right)^{\sqrt{2}} = \left(\sqrt{2}\right)^{\sqrt{2} \sqrt{2}} = \left(\sqrt{2}\right)^2 = 2
    \]
    and $a^b = 2$ is rational.
    \end{itemize}
    In either cases, we can find a pair of irrational numbers $a$ and $b$ such that $a^b$ is rational. Therefore, there are irrational numbers $a$ and $b$ such that $a^b$ is a rational number.
    \end{proof}
\end{solutions*}
